\subsection{Limitations}

The results reported in this study are subject to several limitations, most of which stem from the scarcity of country-level plastic waste data and issues regarding their sources.

First, several series we model are extremely short.
\texttt{YearTRC (Japan)} and \texttt{YearTRC (United States)} each contain nine annual observations.
\texttt{YearPWG (Japan)} contains 26 observations and is evaluated at four forecast origins, with only two observations in each test window.
These windows belong to the same series and provide limited evidence about performance on other datasets.
These small samples limit the complexity of models that can be fit reliably and constrain the strength of conclusions about long-run trends.
With so few test points, an unusual observation or outlier can dominate the reported metrics, and the resulting model rankings are correspondingly unstable.
Therefore, the per-experiment rankings throughout Section~4 should be read as indicative rather than definitive.

Second, some evaluation metrics are unreliable on datasets of this size.
With very few test observations, \(R^2\) is highly sensitive to their values, particularly when the observed values vary little.
For this reason, we omit \(R^2\) from the smallest tasks and rely primarily on error-based metrics such as WAPE.

Third, this study does not account for the organizations and reports from which individual data points are derived.
Because different reports may adopt different methodologies, techniques, and definitions, and because environmental data are inherently noisy, the pooled heterogeneous data exhibit inconsistencies.
We therefore treat this study as a preliminary investigation into data-driven modeling of plastic waste, and we leave the development of models that can accommodate these inconsistencies to future work.

Finally, reproducibility is limited to downstream analyses using the released frozen datasets and code.
The upstream data pipeline is not publicly accessible because access to the PRISM platform and its underlying source records is restricted.
Consequently, external researchers cannot independently repeat the extraction and validation of records from the original heterogeneous documents.
